\section{Introduction}
\label{sec:intro}

\begin{figure}[t]
\centering
\begin{tikzpicture}[
  font=\small,
  layer/.style={draw=black!70, fill=blue!8, rounded corners=1.5pt, minimum width=2.3cm, minimum height=0.42cm, inner sep=1pt, font=\scriptsize},
  mini/.style={draw=black!60, fill=green!10, rounded corners=1pt, minimum width=0.85cm, minimum height=0.2cm, inner sep=0pt},
  tokbox/.style={draw=black!60, fill=orange!12, rounded corners=1pt, minimum width=0.62cm, minimum height=0.42cm, inner sep=1pt, font=\scriptsize},
  clock/.style={-{Latex[length=2.4mm]}, line width=1.1pt},
  thin/.style={-{Latex[length=1.5mm]}, black!55},
  note/.style={font=\scriptsize, align=left},
  ttl/.style={font=\small\bfseries, align=left}
]
% ================= Clock 2: training snapshots =================
\foreach \k/\x/\lab/\op in {0/-7.2/{$\theta_0$}/0.35, 1/-6.0/{$\theta_t$}/0.65, 2/-4.8/{$\theta_{\mathrm{end}}$}/1.0} {
  \begin{scope}[opacity=\op]
    \foreach \i in {0,...,5} { \node[mini] at (\x, 0.42+0.27*\i) {}; }
  \end{scope}
  \node[font=\scriptsize] at (\x, 0.05) {\lab};
}
\draw[clock, green!45!black] (-7.75, -0.3) -- (-4.25, -0.3);
\node[ttl, text=green!40!black, anchor=north west] at (-7.8, -0.45) {Clock 2: training time $t$};
\node[note, anchor=north west, text width=4.0cm] at (-7.8, -0.95) {$\theta_{t+1} = \theta_t - \eta\, g_t(\theta_t)$\\[1pt]edge of stability, catapults and loss spikes, saddle plateaus, grokking};
\draw[thin, dashed] (-4.25, 1.1) -- node[above, font=\scriptsize] {weights $\theta$} (-1.3, 1.1);

% ================= Clock 1: depth =================
\foreach \i/\lab in {0/{embed $x^{0}$}, 1/{block $\ell=1$}, 2/{block $\ell=2$}, 3/{$\cdots$}, 4/{block $\ell=L$}} {
  \node[layer] (L\i) at (0, 0.55*\i) {\lab};
}
\node[draw=black!70, fill=gray!12, rounded corners=1.5pt, minimum width=2.3cm, minimum height=0.42cm, inner sep=1pt, font=\scriptsize] (head) at (0, 2.95) {$\mathrm{softmax}(z/T)$};
\draw[thin] (L4.north) -- (head.south);
\draw[clock, blue!70!black] (1.45, -0.2) -- (1.45, 2.45);
\node[ttl, text=blue!60!black, anchor=north west] at (1.65, 2.55) {Clock 1: layer time $\ell$};
\node[note, anchor=north west, text width=3.9cm] at (1.65, 2.05) {$x^{\ell+1}_i = x^{\ell}_i + f_{\theta}^{\ell}\big(x^{\ell}_i;\, x^{\ell}_{1:n}\big)$\\[1pt]tokens as interacting particles; clustering, rank collapse, order to chaos, dynamical isometry};

% ================= Clock 3: generation =================
\foreach \j/\t in {0/{$y_1$}, 1/{$y_2$}, 2/{$\cdots$}, 3/{$y_{n}$}} {
  \node[tokbox] (T\j) at (-1.05+0.7*\j, -1.0) {\t};
}
\node[tokbox, fill=orange!40] (Tn) at (1.85, -1.0) {$y_{n+1}$};
\draw[thin] (0, -0.77) -- (L0.south);
\draw[clock, orange!80!black, rounded corners=8pt] (head.east) -- (6.0, 2.95) -- (6.0, -1.0) -- (Tn.east);
\node[ttl, text=orange!70!black, anchor=north west] at (-1.4, -1.4) {Clock 3: generation time $n$};
\node[note, anchor=north west, text width=7.2cm] at (-1.4, -1.9) {$y_{n+1} \sim \mathrm{softmax}\big(z(y_{1:n};\theta)/T\big)$: repetition loops as attractors, temperature transitions, drift, collapse};
\end{tikzpicture}
\caption{The three clocks of a transformer language model. \textbf{Clock 1} (blue) iterates the residual update over layers for a fixed input and fixed weights. \textbf{Clock 2} (green) iterates the optimizer over the weights; it hands its current weights to the model. \textbf{Clock 3} (orange) iterates sampling: each new token is appended to the context and the whole stack runs again. One tick of clock~3 contains a full run of clock~1, and clock~2 changes the map that both of them iterate. Typical behavior studied on each clock is listed next to it; Table~\ref{tab:dictionary} maps it onto dynamical-systems concepts.}
\label{fig:clocks}
\end{figure}


A transformer language model is a nonlinear map. It takes a context, a sequence of tokens, and returns a probability distribution over the next token, much as a forced pendulum maps an initial condition and an excitation to a history of displacements.\footnote{Use of AI tools: an AI assistant (Claude, Anthropic) helped search and verify the literature, draft and edit text, and write and run the experiment code, under the author's direction. All references were checked against arXiv or Crossref records and all reported numbers come from the released code. The author takes full responsibility for the content.} The analogy has one important limit. In a mechanical model the parameters, such as the length of a pendulum or the stiffness of a frame under earthquake loading, are fixed or prescribed from outside, whereas a learning model changes its own parameters, and it does so through a second dynamical system, the training algorithm. A language model is therefore used as an iterated map in at least three different ways, and each of these iterations has its own literature, its own vocabulary and its own catalogue of strange behavior.

The first iteration runs over depth. Each residual block adds an update to the hidden state of every token, so a stack of $L$ blocks is a discrete-time dynamical system with the layer index as time \citep{e2017proposal,haber2017stable,lu2017finite,chen2018neural}, and attention couples the tokens into a system of interacting particles \citep{lu2019improving,sander2021sinkformers,geshkovski2023mathematical}. The second iteration runs over optimization steps. Gradient descent is a map on parameter space, and at the step sizes used in practice it sits at the boundary of linear stability \citep{cohen2021gradient}, produces loss spikes \citep{chowdhery2022palm,wortsman2023smallscale} and goes through abrupt transitions such as grokking \citep{power2022grokking} and the formation of induction heads \citep{olsson2022incontext}. The third iteration runs over generated tokens. An autoregressive model feeds its own samples back as input, so decoding is a stochastic dynamical system on the space of contexts, and its best known pathology, the repetition loop \citep{holtzman2019curious,xu2022learning}, has recently been analyzed as an attractor \citep{wang2025unveiling,jones2026chaos}.

We call these three iterations the \emph{three clocks} of a language model, with layer time $\ell$, training time $t$ and generation time $n$ (Figure~\ref{fig:clocks}). The clocks are coupled, since the parameters produced by clock~2 define the map iterated by clock~1, and every tick of clock~3 calls a full run of clock~1. They also run at very different rates, since the weights do not move during generation and each weight update in training averages over many forward passes. Separated time scales of this kind are the standard setting of slow-fast dynamical systems \citep{strogatz2018nonlinear}, and they explain why three literatures could grow up around one object with few citations between them.

The cost of this separation is that well-understood dynamics is rediscovered under new names. Nonlinear dynamics has a mature theory for the behaviors that appear on the three clocks, built over more than a century on pendulums, rotating machinery, ships, lasers and networks of coupled oscillators: stability and bifurcations, synchronization, chaos and Lyapunov exponents, and noise-driven escape from potential wells \citep{strogatz2018nonlinear,kramers1940brownian,hanggi1990reaction,acebron2005kuramoto}. These tools transfer directly. Escape from a repetition loop at sampling temperature $T$ is the escape of a particle from a potential well under noise of intensity $T$, with the gap between the logits of the repeated token and its best competitor in the role of the barrier. The Arrhenius form of the Kramers rate then predicts how the escape probability depends on temperature, and the measured growth of the barrier predicts how long a loop survives (Section~\ref{sec:loops} and Experiment~3). The exchange also runs the other way. Neural networks already forecast spatiotemporally chaotic systems from data \citep{pathak2018model} and learn coordinates in which nonlinear dynamics becomes linear \citep{lusch2017deep}, and a trained transformer is itself a high-dimensional, nonautonomous system that can be run millions of times, observed completely and differentiated exactly.

A bridge between two fields can carry three kinds of traffic. The first is description, in which a behavior seen in one field receives a model, a name and a set of measurable quantities from the other, so that a repetition loop becomes a periodic orbit and a step-to-step oscillation of the loss becomes a period-2 orbit of a map. The second is explanation, in which the tools of the older field account for why the behavior occurs and predict how it changes with a parameter, as Kramers' theory predicts how the rate of escape from a well changes with the noise. The third is solution, in which the explanation points to an intervention that fixes a problem the newer field had not solved. Most of this paper works at the first two levels. Figures~\ref{fig:pair_training} to~\ref{fig:pair_depth} show the description directly, with a mechanical system and a language model drawn in the same representation, and Section~\ref{sec:reduced} explains each correspondence with a reduced-order model, in the way an engineer replaces a sloshing tank or a rolling ship by an equivalent oscillator. The softmax turns out to be a Boltzmann distribution, gradient descent an Euler step on a nonlinear spring, attention on the circle a Kuramoto model, and a repetition loop a well that deepens while it is occupied. At the third level we tested three controllers derived from the escape analysis on GPT-2. They remove exact loops and find the temperature window without a sweep, but none produces more coherent text than the best fixed temperature, because the trap is stored in the context while the controllers act on the noise (\op{12}). Solutions at this level remain the main open task, and the open problems of Section~\ref{sec:open} list the places where we expect them.

In the current work we build this bridge in both directions. We write the three clocks in the language of nonlinear dynamics, record which correspondences are proved, measured or only conjectured, and use the tools of the field to explain behavior that the machine learning literature documents but does not yet explain. Our contributions are the following.
\begin{enumerate}[leftmargin=*,itemsep=2pt]
\item \textbf{A unified framework.} We formalize the three clocks as coupled maps, show textbook physical systems and language models side by side, and derive the reduced-order models behind each correspondence (Section~\ref{sec:primer}) and review the dynamical-systems results for each clock (Sections~\ref{sec:clock1}, \ref{sec:clock2} and~\ref{sec:clock3}), drawing on about 200 papers from machine learning theory, interpretability, optimization, statistical physics and applied mathematics.
\item \textbf{The Three-Clock Dictionary.} Table~\ref{tab:dictionary} maps fifteen concepts from nonlinear dynamics (fixed points, limit cycles, period doubling, Lyapunov exponents, basins, metastability, noise-driven escape, synchronization, Koopman linearization and others) onto observed behavior on each clock, with the evidence level of every entry. An empty cell marks a concept that no study we found has tested on that clock, and each such cell defines a research question.
\item \textbf{Experiments} on open models that run on a laptop CPU (Section~\ref{sec:experiments}). We measure how the per-layer Lyapunov profile of the residual stream changes over training in the public Pythia checkpoints \citep{biderman2023pythia}, we compute a learning-rate bifurcation diagram of a small transformer trained by full-batch gradient descent and show that at the edge of stability several Hessian eigenvalues exceed $2/\eta$ while the curvature along the step stays at $2/\eta$, and we show that escape from repetition loops in GPT-2 and Pythia follows an Arrhenius law with the logit gap as the Kramers barrier \citep{kramers1940brownian}, a barrier that grows while the model stays in the loop and then saturates. We prove that logarithmic growth would give a sharp locking temperature, show that the measured growth saturates instead, so that the locking temperature falls with the length of the generation and rises with model size from 70M to 1B parameters, and map the narrow window of temperatures between repetitive and incoherent text. In a small reasoning model the same window appears, centered on the temperature its developers recommend, but its lower edge is set by recurrent rethinking, and no exact loop forms.
\item \textbf{Twelve open problems} (Section~\ref{sec:open}), each stated with a first test.
\end{enumerate}

Several papers develop one of these views in depth, and we rely on them throughout. \citet{geshkovski2023mathematical} and \citet{castin2025unified} develop the mathematics of attention as an interacting particle system over depth. \citet{cohen2021gradient} and the work that followed analyze gradient descent at the edge of stability. \citet{krishnamurthy2026stochastic} give a stochastic dynamical-systems account of autoregressive generation and of networks of interacting language models, and \citet{jones2026chaos} studies chaos and bifurcations in a deterministic feedback version of generation. To our knowledge, no earlier work treats the three clocks side by side in one vocabulary, with an explicit record of which correspondences are proved, which are measured and which are conjectures.

We use the language of dynamical systems in its technical sense wherever a result supports it, and we flag analogies as analogies. For example, we call a loss spike a bifurcation only when a parameter can be shown to cross a stability threshold. The paper is written for machine learning researchers who want a compact entry into dynamical-systems tools, and for dynamicists who want to know which of their tools are already in use on language models. Section~\ref{sec:primer} is a primer for the first group, and the dictionary in Section~\ref{sec:dictionary} is intended for both.

\begin{center}
\fbox{\begin{minipage}{0.94\linewidth}
\textbf{Practical takeaways.}
\begin{itemize}[leftmargin=*,itemsep=1pt,topsep=2pt]
\item \emph{Generation.} In GPT-2 and Pythia, coherent text appears only in a narrow band of temperatures between repetition and incoherence, and for some prompts in none. A loop is not locked by a fixed temperature: the temperature needed to leave it within a given budget falls like $1/\ln N$ with the length $N$ of the generation, and larger models need a higher temperature to leave the same kind of loop (Section~\ref{sec:exp3}). In Qwen3-0.6B on GSM8K, accuracy within a fixed budget peaks at the vendor's $T=0.6$, and below it the model rethinks the same step until the budget runs out.
\item \emph{Anti-repetition.} Raising the temperature when a loop starts, rolling the text back to the start of the loop, or steering the temperature by feedback all remove exact loops, but none beats the best fixed temperature on coherent text, because the attractor is stored in the context, not in the sampler (\op{12}).
\item \emph{Training.} At the edge of stability several Hessian eigenvalues sit 5\% to 14\% above $2/\eta$ while the curvature along the actual step equals $2/\eta$, so $\eta\lambda_{\max}>2$ alone does not signal divergence. A power iteration can underestimate $\lambda_{\max}$ there, because the top eigenvalues form a cluster that reorders at every step (Section~\ref{sec:exp2}).
\item \emph{Depth.} Training makes each layer stretch perturbations more, while the stretched directions of consecutive layers become less aligned (Section~\ref{sec:exp1}).
\item \emph{Tools.} Per-layer Lyapunov exponents, learning-rate bifurcation diagrams and exact loop-escape probabilities from a single greedy pass all run on a laptop CPU, and the code is released.
\end{itemize}
\end{minipage}}
\end{center}

